\chapter{Operatore SELECT - Interrogazione di un database}
\section{A cosa serve?}
È la istruzione utilizzata per effettuare interrogazioni al database

\section{Sintassi}
\begin{lstlisting}[language=sql_generic]
SELECT ListaAttributi
FROM ListaTabelle
[WHERE Condizione]
\end{lstlisting}

\section{Esempio}
\begin{lstlisting}[language=sql_generic]
SELECT Nome, Cognome, Stipendio
FROM Impiegato
WHERE Mansione = 'Ingegnere'
\end{lstlisting}

\section{Istruzione AS}
\subsection{A cosa serve?}
Viene utilizzato per fare ridenominazioni su attributi e anche su relazioni

\subsection{Sintassi}
\begin{lstlisting}[language=sql_generic]
SELECT AttrEspr [ [ AS ] Alias ] {, AttrEspr [ [ AS ] Alias ] }
FROM Tabella [ [ AS ] Alias ] {, Tabella [ [ AS ] Alias ] }
[ WHERE Condizione ]
\end{lstlisting}

\subsection{Esempio}
\begin{lstlisting}[language=sql_generic]
SELECT I.Nome, D.Citta AS citta_lavoro
FROM Impiegato AS I, Dip AS D
WHERE I.Dipart = D.Nome;
\end{lstlisting}

\section{Clausola FROM}
Specifica la/le tabelle/e su cui operare per determinare il risultato

\subsection{Sintassi}
\begin{lstlisting}[language=sql_generic]
SELECT ListaAttributi
FROM Tabella1 [ [ AS ] Alias1 ] {, TabellaN [ [ AS ] AliasN ] }
\end{lstlisting}

\subsection{Notazione punto}
La clausola \textbf{.} permette di specificare a che relazione appartiene un attributo (o anche una tabella a un database)

\subsection{Esempio}
\begin{lstlisting}[language=sql_generic]
SELECT I.Nome, D.NomeDip
FROM Impiegato AS I, Dip AS D
WHERE I.Dipart = D.CodDip;
\end{lstlisting}

\section{Clausola WHERE}
\subsection{A cosa serve?}
È una espressione booleana che specifica una o più condindizioni di selezione

\subsection{Operatori di confronto}
\begin{itemize}
    \item =
    \item <>
    \item <= =>
    \item LIKE
    \item BETWEEN
    \item IS NULL
    \item IS NOT NULL
\end{itemize}

\subsection{Operatore LIKE}
L'operatore LIKE permette di esprimere dei "pattern" su stringhe di caratteri (regex) mediante le seguente notazione:
\begin{itemize}
    \item \_ indica un singolo carattere arbitrario
    \item \% indica una strigna di caratteri arbitraria (anche lunga 0)
\end{itemize}

\subsection{Operatore BETWEEN}
L'operatore BETWEEN permette di esprimere condizioni di appartenenza a un intervallo

\subsection{Operatore IN}
L'operatore IN permette di esprimere condizioni di appartenenza a un insieme

\subsection{Predicati semplici}
Sono espressioni di valutazione su attributi mediante operatori di confronto

\subsection{Precedenza di valutazione tra operatori booleani}
\begin{enumerate}
    \item NOT
    \item AND
    \item OR
\end{enumerate}

\subsection{Sintassi}
\begin{lstlisting}[language=sql_generic]
WHERE [NOT] PredicatoSemplice {< AND | OR > [NOT] PredicatoSemplice}
\end{lstlisting}

\subsection{Esempio generico}
\begin{lstlisting}[language=sql_generic]
SELECT Nome, Cognome, Mansione
FROM Impiegato
WHERE ufficio = '10';
\end{lstlisting}

\subsection{Esempio con LIKE}
\begin{figure}[htpb]
\begin{lstlisting}[language=sql_generic]
SELECT Nome, Cognome 
FROM Impiegato 
WHERE Cognome LIKE '%i';
\end{lstlisting}
\caption{Restituira tutti i cognomi che finiscono con la lettera i}
\end{figure}
\begin{figure}[htpb]
    \begin{lstlisting}[language=sql_generic]
SELECT Nome, Cognome 
FROM Impiegato 
WHERE Cognome LIKE 'R_ssi';
    \end{lstlisting}
    \caption{Restituirà "Rossi" (dove \_ sostituisce la 'o'), ma non "Rossini" (perché \_ impone un solo carattere)}
\end{figure}

\subsection{Esempio con BETWEEN}
\begin{lstlisting}[language=sql_generic]
SELECT Nome, Stipendio 
FROM Impiegato 
WHERE Stipendio BETWEEN 38 AND 60;
\end{lstlisting}

\subsection{Esempio con IN}
\begin{lstlisting}[language=sql_generic]
SELECT Cognome, ufficio 
FROM Impiegato 
WHERE ufficio IN ('10', '20');
\end{lstlisting}

\section{Clausola DISTINCT}
\subsection{A cosa serve?}
In SQL, le tabelle prodotte dalle interrogazioni possono contenere più righe identiche tra loro.

I duplicati possono essere rimossi usando la parola chiave DISTINCT

\subsection{Sintassi}
\begin{lstlisting}[language=sql_generic]
SELECT DISTINCT ListaAttributi
FROM ListaTabelle
[ WHERE Condizione ]
\end{lstlisting}

\subsection{Esempio}
\begin{lstlisting}[language=sql_generic]
SELECT DISTINCT Cognome, Citta 
FROM Impiegato 
WHERE Stipendio > 40;
\end{lstlisting}

\section{Espressioni}
\subsection{A cosa servono?}
Una espressione usata nella clausola SELECT dà luogo ad una nuova "colonna" che non corrisponde a nessuna colonna della tabella su cui si effettua la selezione.

Questa nuova colonna è detta \underline{colonna virtuale}.

Le colonne virtuali non sono fisicamente memorizzate, ma esistono solo come risultato delle interrogazioni.

Anche alle colonne virtuali è possibile assengare un alias (con il costrutto AS)

\subsection{Funzioni predefinite}
SQL mette a disposizione alcune funzioni predefinite:
\begin{itemize}
    \item \textbf{Stringhe:} UPPER, UCASE, LENGTH, \dots
    \item \textbf{Date / Intervalli:} +, -, DATE,DAYOFWEEK, \dots
    \item \textbf{Matematiche:} * , + , - , / , TAN,SQRT, SIN, \dots
    \item \textbf{Informazioni di sistema:} USER, CURRENT\_DATE, \dots
\end{itemize}

\subsection{Funzioni per stringhe}
Gli operatori più comuni per le stringhe sono:
\begin{itemize}
    \item \textbf{||:} Operatore di concatenazione
    \item \textbf{UPPER, UCASE:} Trasforma la stringa in caratteri maiuscoli
    \item \textbf{LOWER, LCASE:} Trasforma la stringa in caratteri minuscoli
\end{itemize}

\subsection{Esempio di utilizzo di funzioni aritmetiche}
\begin{lstlisting}[language=sql_generic]
SELECT Nome, Stipendio, Premioprod, (Stipendio + Premioprod) AS stipendio_totale
FROM Impiegato
WHERE Mansione = 'Ingegnere';
\end{lstlisting}
\subsection{Esempio di utilizzo di funzioni per stringhe}
\begin{lstlisting}[language=sql_generic]
SELECT Nome, Cognome, citta, mansione
FROM Impiegato
WHERE UCASE(mansione) = 'INGEGNERE';
\end{lstlisting}